\subsection{可分解李代数的特征刻画}

每个可分解李代数作为向量空间（更不必说作为李代数）由它的半单或幂零元素所成的集生成。反之：

定理 1. 设 g 是 \(\mathfrak{gl}(V)\) 的一个李子代数，并设 X 是 g 的一个子集，它把 g 生成为 k 上的李代数。若 X 的每个元素都是半单的或幂零的，则 g 可分解。

\begin{enumerate}
\def\labelenumi{\alph{enumi})}
\tightlist
\item
  \(\mathfrak{g}\) 交换。
\end{enumerate}

g 的半单（相应地，幂零）元素构成一个向量子空间 \(g_{s}\)（相应地，\(g_{n}\)）。假定等价于 \(g = g_{s} \oplus g_{n}\)，由此得到 g 可分解这一事实。

\begin{enumerate}
\def\labelenumi{\alph{enumi})}
\setcounter{enumi}{1}
\tightlist
\item
  \(\mathfrak{g}\) 是约化的。
\end{enumerate}

则 \(\mathfrak{g} = \mathfrak{g}' \times \mathfrak{c}\)，其中 \(\mathfrak{g}'\) 半单，\(\mathfrak{c}\) 交换。由命题 2，\(\mathfrak{g}'\) 可分解。设 \(x = a + b \in \mathfrak{g}\)，其中 \(a \in \mathfrak{g}'\)，\(b \in \mathfrak{c}\)。设 \(a_s, a_n, b_s, b_n\) 是 \(a, b\) 的半单分量与幂零分量。由于 \(a_s, a_n, b_s, b_n\) 相互交换，\(x\) 的半单分量与幂零分量是 \(a_s + b_s, a_n + b_n\)。而现在 \(a_s, a_n \in \mathfrak{g}'\)。若 \(x\) 半单，则 \(x = a_s + b_s\)；由于 \(a_s \in \mathfrak{g}'\)，有 \(b_s \in \mathfrak{g}\)，故 \(b_s \in \mathfrak{c}\)，因为 \(b_s\) 与 \(\mathfrak{g}\) 交换；因而 \(a = a_s\) 且 \(b = b_s\)。类似地，若 \(x\) 幂零，则 \(a = a_n\) 且 \(b = b_n\)。由此推出，X 的元素在 \(\mathfrak{c}\) 上的投影是半单的或幂零的；由 \(a)\)，这蕴含 \(\mathfrak{c}\) 可分解。保留前面的记号，但不加于 \(x\) 的假定，我们现在有 \(b_s, b_n \in \mathfrak{c}\)，于是 \(a_s + b_s, a_n + b_n \in \mathfrak{g}\)，这就证明了定理在这种情形下成立。

\begin{enumerate}
\def\labelenumi{\alph{enumi})}
\setcounter{enumi}{2}
\tightlist
\item
  一般情形。
\end{enumerate}

我们设定理对维数 \(< \dim \mathfrak{g}\) 的李代数已证明，并对 \(\mathfrak{g}\) 证明它。

设 \(\mathfrak{n}\) 是 \(\mathfrak{g}\) 的恒等表示的最大幂零性理想。若 \(\mathfrak{n} = 0\)，则 \(\mathfrak{g}\) 有一个单射半单表示，从而是约化的。设 \(\mathfrak{n} \neq 0\)。设 \(\mathfrak{p}\) 是 \(\mathfrak{n}\) 在 \(\mathfrak{gl}(\mathrm{V})\) 中的正规化子。存在满足命题 8 条件的 E、\(\rho, \mathrm{F}\)。由于 \(\mathfrak{g} \subset \mathfrak{p}\)，\(\rho(\mathfrak{g})\) 使 F 保持稳定；设 \(\rho_0\) 是表示 \(u \mapsto \rho(u)|\mathrm{F}\)（\(\mathfrak{g}\) 在 F 上的表示）；我们有 \(\mathfrak{n} = \mathrm{Ker}\rho_0\)。在 \(\rho\) 下，\(\mathfrak{gl}(\mathrm{V})\) 的每个半单（相应地，幂零）元素的像是半单的（相应地，幂零的）（命题 2）。于是代数 \(\rho_0(\mathfrak{g})\) 由它的半单元素与幂零元素生成。由归纳假设，\(\rho_0(\mathfrak{g})\) 可分解。

设 \(x \in \mathfrak{g}\)，并设 \(x_s, x_n\) 是它的半单分量与幂零分量。由命题 2，\(\rho(x)\) 的半单分量与幂零分量是 \(\rho(x_s), \rho(x_n)\)。由于 \(\rho_0(\mathfrak{g})\) 可分解，存在 \(y, z \in \mathfrak{g}\)，使得

\[
\rho_ {0} (y) = \rho (x _ {s}) | \mathrm{F}, \rho_ {0} (z) = \rho (x _ {n}) | \mathrm{F}.
\]

则 \(x_{s}\in y + \mathfrak{n},x_{n}\in z + \mathfrak{n}\)，故 \(x_{s},x_{n}\in \mathfrak{g}\)

证毕。

推论 1. \(\mathfrak{gl}(\mathrm{V})\) 的每个由它的可分解子代数生成的子代数都是可分解的。

这是明显的。

推论 2. 设 \(\mathfrak{g}\) 是 \(\mathfrak{gl}(\mathrm{V})\) 的一个李子代数。则 \([\mathfrak{g},\mathfrak{g}]\) 可分解。

设 \(\mathfrak{r}\) 是 \(\mathfrak{g}\) 的根基，\(\mathfrak{s}\) 是 \(\mathfrak{g}\) 的一个 Levi 子代数（第一章，§6，第 8 节）。则

\[
[ \mathfrak {g}, \mathfrak {g} ] = [ \mathfrak {s}, \mathfrak {s} ] + [ \mathfrak {s}, \mathfrak {r} ] + [ \mathfrak {r}, \mathfrak {r} ] = \mathfrak {s} + [ \mathfrak {g}, \mathfrak {r} ].
\]

代数 \([g,r]\) 可分解，因为它的所有元素都是幂零的（第一章，§5，第 3 节）。另一方面，s 可分解（命题 2）。由此推出 \([g,g]\) 可分解（推论 1）。

推论 3. 设 \(\mathfrak{g}\) 是 \(\mathfrak{gl}(\mathrm{V})\) 的一个李子代数，并设 \(\mathrm{X}\) 是 \(\mathfrak{g}\) 的一个生成 \(\mathfrak{g}\) 的子集（作为 \(k\) 上的李代数）。

\begin{enumerate}
\def\labelenumi{(\roman{enumi})}
\item
  可分解包络 \(e(\mathfrak{g})\)（\(\mathfrak{g}\) 的可分解包络）由 \(X\) 的元素的半单分量与幂零分量生成。
\item
  若 \(k'\) 是 \(k\) 的一个扩张，则 \(e(\mathfrak{g} \otimes_k k') = e(\mathfrak{g}) \otimes_k k'\)；且 \(\mathfrak{g}\) 可分解当且仅当 \(\mathfrak{g} \otimes_k k'\) 可分解。
\end{enumerate}

设 \(\tilde{\mathfrak{g}}\) 是 \(\mathfrak{gl}(\mathrm{V})\) 的由 X 的元素的半单分量与幂零分量生成的子代数。则 \(\mathfrak{g} \subset \tilde{\mathfrak{g}} \subset e(\mathfrak{g})\)；由定理 1，\(\tilde{\mathfrak{g}}\) 可分解，故 \(\tilde{\mathfrak{g}} = e(\mathfrak{g})\)，这就证明了 (i)。断言 (ii) 由此得出，因为 X 生成 \(k'\) -代数 \(\mathfrak{g} \otimes_k k'\)。

推论 4. 设 g 是 \(\mathfrak{gl}(\mathrm{V})\) 的一个可分解李子代数。设 T 是由半单元素组成的 g 的交换子代数所成的集（参见命题 6）。T 的极大元素都有相同的维数。

设 \(k'\) 是 \(k\) 的一个代数闭扩张，且 \(V' = V \otimes_k k', \mathfrak{g}' = \mathfrak{g} \otimes_k k'\)。设 \(t_1, t_2\) 是 \(\mathcal{T}\) 的极大元素，\(t_i' = t_i \otimes_k k'\)，\(\mathfrak{h}_i\) 是 \(t_i\) 在 \(\mathfrak{g}\) 中的交换子，\(\mathfrak{h}_i' = \mathfrak{h}_i \otimes_k k'\)。则 \(\mathfrak{h}_i\) 是 \(\mathfrak{g}\) 的一个嘉当子代数（命题 6），故 \(\mathfrak{h}_i'\) 是 \(\mathfrak{g}'\) 的一个嘉当子代数。于是 \(\mathfrak{h}_i = t_i \times \mathfrak{n}_V(\mathfrak{h}_i)\)，因而 \(\mathfrak{h}_i' = t_i' \times \mathfrak{n}_{V'}(\mathfrak{h}_i')\)，于是 \(t_i'\) 是 \(\mathfrak{h}_i'\) 的半单元素所成的集。由于 \(\mathfrak{g}'\) 可分解（推论 3），\(t_1'\) 与 \(t_2'\) 在 \(\text{Aut}_e(\mathfrak{g}')\) 下共轭（命题 6），故 \(\dim t_1 = \dim t_2\)。

定理 2. 设 \(\mathfrak{g}\) 是 \(\mathfrak{gl}(\mathrm{V})\) 的一个李子代数。下列条件等价：

\begin{enumerate}
\def\labelenumi{(\roman{enumi})}
\item
  \(\mathfrak{g}\) 可分解；
\item
  \(\mathfrak{g}\) 的每个嘉当子代数都可分解；
\item
  \(\mathfrak{g}\) 有一个可分解的嘉当子代数；
\item
  \(\mathfrak{g}\) 的根基可分解。
\item
  \(\Longrightarrow\) (ii)：这由命题 3 的推论 2 得出。
\item
  \(\Longrightarrow\) (i)：这由定理 1 的推论 1 得出，因为 \(\mathfrak{g}\) 由它的嘉当子代数生成（§2，第 3 节，定理 1 的推论 3）。
\item
  \(\Longrightarrow\) (iii)：这是明显的。
\item
  \(\Longrightarrow\) (ii)：由定理 1 的推论 3，可以设 k 代数闭。于是 g 的嘉当子代数在 g 的初等自同构下共轭（§3，第 2 节，定理 1）；鉴于 §3 第 1 节的注记 1，由此推出，若其中一个是可分解的，则它们全都是可分解的。
\item
  \(\Longrightarrow\) (iv)：这由命题 4 的推论 2 得出。
\item
  \(\Longrightarrow\) (i)：设 g 的根基 r 可分解。设 s 是 g 的一个 Levi 子代数；它可分解（命题 2）。因而 \(g = s + r\) 可分解（定理 1 的推论 1）。
\end{enumerate}

